\documentclass[a4,10pt]{article}
%%%%%%% --------------------------------------------------------------------------------------
%%%%%%% PREAMBLE
%%%%%%% --------------------------------------------------------------------------------------
%%%%%%% Compile with XeLaTeX or LuaLaTeX (fontspec requires it).
%%%%%%% Needs SourceSansPro-*.otf on the font path and Publications.bib alongside.
%%%%%%%
%%%%%%% CHANGED 2026-07: Microsoft AI (April 2026--present) added as current role;
%%%%%%% Meta updated to March 2025--April 2026.
%%%%%%% --------------------------------------------------------------------------------------
\usepackage{latexsym}
\usepackage[empty]{fullpage}
\usepackage{titlesec}
\usepackage{marvosym}
\usepackage[usenames,dvipsnames]{color}
\usepackage{verbatim}
\usepackage[hidelinks]{hyperref}
\usepackage{fancyhdr}
\usepackage{multicol}
\usepackage{csquotes}
\usepackage{tabularx}
\usepackage[11pt]{moresize}
\usepackage{setspace}
\usepackage{fontspec}
\usepackage[inline]{enumitem}
\usepackage{array}
\newcolumntype{P}[1]{>{\centering\arraybackslash}p{#1}}
\usepackage{anyfontsize}
%%%% Set Margins
\usepackage[margin=1cm, top=1cm]{geometry}
%%%% Set Fonts
\setmainfont[
  BoldFont=SourceSansPro-Semibold.otf,
  ItalicFont=SourceSansPro-RegularIt.otf
]{SourceSansPro-Regular.otf}
\setsansfont{SourceSansPro-Semibold.otf}
%%%% CJK family for the Chinese form of the name (XeLaTeX/LuaLaTeX only).
%%%% If "PingFang SC" is unavailable, swap for another installed CJK font
%%%% (e.g. "Songti SC", "Noto Serif CJK SC") or comment out both this line
%%%% and the \cjkname use in the header below.
\newfontfamily\cjkname{PingFang SC}
%%%% Page Style
\pagestyle{fancy}
\fancyhf{}
\fancyfoot{}
\renewcommand{\headrulewidth}{0pt}
\renewcommand{\footrulewidth}{0pt}
%%%% URL Style
\urlstyle{same}
%%%% Indentation / Layout
\raggedbottom
\raggedright
\setlength{\tabcolsep}{0in}
%%%% Secondary Color
\definecolor{UI_blue}{RGB}{32, 64, 151}
%%%% Bibliography / Name Highlighting
\usepackage[style=nature, maxbibnames=3]{biblatex}
\addbibresource{Publications.bib}
\renewcommand*{\mkbibnamegiven}[1]{%
  \ifitemannotation{highlight}{\textbf{#1}}{#1}}
\renewcommand*{\mkbibnamefamily}[1]{%
  \ifitemannotation{highlight}{\textbf{#1}}{#1}}
%%%% Section Formatting
\titleformat{\section}{
  \color{UI_blue}
  \scshape
  \raggedright
  \large
}{}{0em}{}[\vspace{-8pt}\hrulefill\vspace{-4pt}]
\titlespacing*{\subsection}{1pt}{0.8em}{0.5em}
\titleformat{\subsection}{\vspace{-0.05cm}\bfseries\fontsize{11pt}{2cm}}{}{0em}{}[\vspace{-0.15cm}]
%%%% Subtext Formatting
\newcommand{\subtext}[1]{#1\par\vspace{-0.2cm}}
%%%% Item Spacing
\setlist[itemize]{align=parleft,left=0pt..1em}
%%%% New Itemize "Zitemize" Environment
\newenvironment{zitemize}{
  \begin{itemize}\itemsep0pt\parskip0pt\parsep1pt}
  {\end{itemize}\vspace{-0.4cm}}
%%%% Skills Bold Formatting
\newcommand{\hskills}[1]{\textbf{\bfseries #1}}
%%%% Icons
\usepackage{fontawesome5}
\usepackage{academicons}
\usepackage{xcolor} % required by academicons

\begin{document}

% ==========================================================
% HEADER
% ==========================================================
\begin{center}
{\Huge \textbf{Hans William Alexander Hanley}}\ {\large\color{UI_blue}\cjkname{}汉斯·汉隶}\\[2pt]
{\large \color{UI_blue} Semantic Modeling • Natural Language Processing • LLM Safety • Multilingual Representation Learning}\\[6pt]

\begin{tabular}{c c c}
\faEnvelope\ \href{mailto:hhanley@cs.stanford.edu}{hhanley@cs.stanford.edu} &
\faGlobe\ \href{https://www.hanshanley.com}{hanshanley.com} &
\faPenNib\ \href{https://www.themarginoferror.com}{themarginoferror.com} \\
\faLinkedin\ \href{https://www.linkedin.com/in/hans-hanley-0694a180}{hans-hanley-0694a180} &
\faGithub\ \href{https://github.com/hanshanley}{github.com/hanshanley} &
\aiGoogleScholar\ \href{https://scholar.google.com/citations?user=ewdWfOoAAAAJ&hl=en&oi=ao}{Google Scholar}
\end{tabular}
\end{center}

% ==========================================================
% SUMMARY
% ==========================================================
I build scalable methods for interpreting and evaluating large language models, including multilingual representation-learning systems, semantic modeling pipelines, and LLM-based evaluators and agents for high-fidelity reasoning and content analysis. My PhD research uncovered latent structure in massive text ecosystems, aligned LLM behavior across languages and domains, and supported safer, more reliable deployment of model-driven systems.

% ==========================================================
% EDUCATION
% ==========================================================
\vspace{-1.5em}
\section{Education}
\vspace{-0.5em}

\subsection*{Stanford University \hfill PhD, Computer Science, 2025 \,|\, MS, Computer Science, 2024}
\subtext{GPA: 4.02; Advisor: Zakir Durumeric; Meta Research PhD Fellow; NSF GRFP; Stanford EDGE Doctoral Fellow}

\subsection*{University of Oxford \hfill MSc, Statistical Science, 2020 \,|\, MSc, Computer Science, 2019}
\subtext{Distinction (Highest Honors); Daniel M.~Sachs Class of 1960 Graduating Scholar at Worcester College}

\subsection*{Princeton University \hfill BSE, Electrical Engineering, 2018}
\subtext{Highest Honors, GPA: 3.982; Information Security \& Privacy Concentration; Phi Beta Kappa; Tau Beta Pi; Sigma Xi}

% ==========================================================
% EXPERIENCE
% ==========================================================
\vspace{-0.8em}
\section{Professional Experience}
\vspace{-0.5em}

\subsection*{Microsoft AI, Member of Technical Staff \hfill April 2026--Present}
\begin{zitemize}
\item Research and build large-scale language-model systems for AI safety and truthfulness, spanning multilingual representation learning, semantic modeling, and evaluation of model-generated content.
\item Develop LLM-based evaluators and measurement pipelines to assess factuality, reasoning quality, and safety and integrity risks in deployed model-driven products.
\end{zitemize}

\vspace{-0.25em}
\subsection*{Meta, Research Scientist, CAS / Central Applied Sciences \hfill March 2025--April 2026}
\begin{zitemize}
\item Researched, built, and deployed large-scale multilingual and multimodal LLM pipelines to model content semantics, identify creator verticals, and understand user interaction behavior across Instagram (over 3 billion users) for downstream recommender systems.
\item Developed LLM-based evaluation methods and production metrics to assess political content quality, reply/comment impact, safety and integrity risks, and downstream effects on user experience.
\end{zitemize}

\vspace{-0.25em}
\subsection*{Google Research, Health AI Student Researcher \hfill Summer 2024}
\begin{zitemize}
\item Designed LLM-based agents for structured querying and high-fidelity reasoning over large clinical datasets; implemented verification layers to detect reasoning failures under distribution shifts.
\item Built evaluation frameworks measuring accuracy, safety, and adversarial robustness for LLM-mediated medical data interactions.
\end{zitemize}

\vspace{-0.25em}
\subsection*{Digital Forensics Research Lab, Atlantic Council — Research Intern \hfill Summer 2021}
\begin{zitemize}
\item Investigated global misinformation and influence operations; authored seven reports analyzing cross-platform inauthentic behavior, foreign influence operations, and propaganda narratives (e.g., COVID-19 conspiracy networks, Iran/Russia/China operations, VPN adoption in Nigeria).
\end{zitemize}

% ==========================================================
% RESEARCH HIGHLIGHTS & PUBLICATIONS
% ==========================================================
\vspace{-1em}
\section{Research Highlights \& Selected Publications}
\vspace{-0.5em}

\subsection*{\href{https://www.usenix.org/system/files/conference/usenixsecurity25/sec25cycle1-prepub-487-hanley.pdf}{Tracking the Takes and Trajectories of News Narratives} \hfill USENIX Security 2025}
\begin{zitemize}
\item Built a large-scale narrative-tracking system combining specially fine-tuned (e.g., LoRA) encoder-LLM semantic similarity models and zero-shot stance detection to track 146K narratives across 4K+ news sites and understand narrative pathways (e.g., anti-vaccine and anti-Ukraine influence campaigns).
\end{zitemize}

\vspace{-0.25em}
\subsection*{\href{https://www.pnas.org/doi/10.1073/pnas.2420607122}{Across the Firewall: Foreign Media’s Influence on Chinese Social Media} \hfill PNAS 2025}
\begin{zitemize}
\item Quantified foreign-media penetration into Chinese platforms using cross-lingual alignment pipelines to detect narrative borrowing on the Russo--Ukrainian War.
\end{zitemize}

\vspace{-0.25em}
\subsection*{\href{https://aclanthology.org/2025.acl-long.124/}{Hierarchical Level-Wise News Article Clustering via Multilingual Matryoshka Embeddings} \hfill ACL 2025}
\begin{zitemize}
\item Proposed multilingual Matryoshka embeddings enabling controllable semantic granularity (coarse-to-fine); achieved SOTA on SemEval 2022 Task 8 and built scalable hierarchical clustering pipelines surfacing stories, narratives, and macro-themes.
\end{zitemize}

\vspace{-0.25em}
\subsection*{\href{https://aclanthology.org/2023.emnlp-main.694/}{TATA: Stance Detection via Topic-Agnostic and Topic-Aware Embeddings} \hfill EMNLP 2023}
\begin{zitemize}
\item Developed contrastive embedding models for topic-agnostic and topic-aware stance detection; achieved SOTA zero-shot performance (0.771 F1 on VAST) and improved cross-topic generalization in multi-domain news corpora.
\end{zitemize}

% ==========================================================
% SKILLS
% ==========================================================
\vspace{-1em}
\section{Skills}
\vspace{-0.5em}

\begin{tabular}{p{7em} p{48em}}
\hskills{Programming} & Python, Java, C, C++, C\#, Go, SQL, R, MATLAB; Git; PyTorch, TensorFlow, Hugging Face \\[4pt]
\hskills{ML/LLM} & Transformers; LLM fine-tuning; multilingual modeling; interpretability; feature attribution; contrastive learning; LLM-based evaluation; safety \& robustness analysis \\[4pt]
\hskills{Languages} & English (native); Mandarin Chinese (HSK Level 5) \\
\end{tabular}

% ==========================================================
% AWARDS
% ==========================================================
\vspace{-1em}
\section{Selected Awards \& Honors}
\vspace{-0.5em}

\begin{zitemize}
\item \href{https://medium.com/acm-cscw/announcing-the-best-of-cscw-2025-a95517e67ba3}{CSCW 2025 Best Paper Honorable Mention} (Top 4\%).
\item \href{https://datascience.stanford.edu/news/shining-light-emerging-talent-rising-stars-data-science-2024-recap}{Rising Stars in Data Science (Stanford/UCSD/UChicago)};
\href{https://research.facebook.com/fellows/hanley-hans-william-alexander/}{Meta Research PhD Fellow};
\href{https://ece.princeton.edu/news/seven-current-and-former-students-awarded-nsf-graduate-research-fellowships}{NSF Graduate Research Fellowship (GRFP)};
\href{https://vpge.stanford.edu/fellowships-funding/current-vpge-fellows/all-2020}{Stanford EDGE Doctoral Fellow};
\href{https://www.princeton.edu/news/2017/12/11/princeton-seniors-hanley-and-silver-oxford-graduate-student-barnett-named-sachs}{University of Oxford Daniel M.\ Sachs Class of 1960 Graduating Scholar}.
\item \href{https://ece.princeton.edu/news/undergraduates-honored-class-day-2018}{Charles Ira Young Memorial Tablet \& Medal};
\href{https://engineering.princeton.edu/news/2018/06/05/class-day-awards-celebrate-graduates-contributions-research-and-service}{SEAS Hayes--Palmer Prize};
\href{https://www.princeton.edu/news/2016/09/11/princeton-students-honored-opening-exercises}{George B.\ Wood Legacy Prize}.
\end{zitemize}

\end{document}
